% id: 119-07
% book: 베이직쎈 확률과 통계 · 06 이항분포와 정규분포
% section: 유형 075 이항분포와 정규분포의 관계의 활용
% figure: tikz:fig-119-07
% answer: 
% answer_source: 
% variant_level: 0 (원본 전사)
% 이 파일은 build.mjs 가 items.json 에서 만든다. 고칠 때는 items.json 을 고친다.
\smash{\raisebox{1.5mm}{\dmkichul{베이직쎈 확률과 통계 119쪽 07번}}}
\noindent\begin{minipage}[t]{0.58\linewidth}\vspace{0pt}\raggedright 좌석이 348개인 어느 비행기를 예약했다가 취소할 확률은 $10\mkern3mu \%$라 한다. 이 비행기를 예약한 400명의 고객 중 예약을 취소하는 고객의 수를 $X$라 할 때, 다음에 답하시오. \dmcondfill{예약을 취소하는 시행은 독립시행이다.}{}\end{minipage}\hfill\begin{minipage}[t]{0.4\linewidth}\vspace{0pt}\centering \resizebox{\ifdim\width>\linewidth\linewidth\else\width\fi}{!}{% 119-07(인쇄 119쪽) — 119-07 소문항 (3) 오른쪽에 놓인 표준정규분포표($z$ / $\mathrm{P}(0\le Z\le z)$).
% 스캔 표라 크롭하지 않고 tabular 로 다시 조판했다(칸 값은 원문 그대로 · 원문에 없는 행은 넣지 않는다).
{\setlength{\tabcolsep}{5pt}\renewcommand{\arraystretch}{1.25}%
\begin{tabular}{|c|c|}
\hline
$z$ & $\mathrm{P}(0\le Z\le z)$ \\
\hline
$1.0$ & $0.3413$ \\
\hline
$1.5$ & $0.4332$ \\
\hline
$2.0$ & $0.4772$ \\
\hline
\end{tabular}}
}\end{minipage}\par
\dmsubprob{1} 확률변수 $X$가 따르는 이항분포를 $\mathrm{B}(n{,}\mkern3mu\mathopen{}p)$ 꼴로 나타내시오.
\dmsubprob{2} 확률변수 $X$가 근사적으로 따르는 정규분포를 $\mathrm{N}(m{,}\mkern3mu\mathopen{}\sigma^2)$ 꼴로 나타내시오.
\dmsubprob{3} 오른쪽 표준정규분포표를 이용하여 이 비행기를 타러 온 모든 고객이 비행기를 탈 확률을 구하시오.
